\scanpage{45}\mspage{$\underline{\overline{\mathrm{XVII}}}$}
$[s,s']\subset A(\Lambda^{1})$ (resp. $A(\Lambda^{2})$) ce qui est absurde \unsure{puìsque|puisque}
\[
{]s,s'[}\subset\mathbb{R}^{2}-\Lambda^{1}\quad(\text{resp. }\mathbb{R}^{2}-\Lambda^{2}).
\]

\uline{Lemme} $\underline{\overline{\mathrm{XVIII}}}$. $\Lambda^{2}-\{s,s'\}\deleted{\unsure{\supset}}\subset\mathrm{ext}\,\Lambda^{1}$ et $\Lambda^{1}-\{s,s'\}\subset\mathrm{ext}\,\Lambda^{2}$.

\uline{Démonstration} : Soit $D_s$ un disque de centre $s$ satisfaisant b) de la def. 1 par rapport a $\Lambda^{1}$ \deleted{et}, à $\Lambda^{2}$ et \={a} $\Lambda$.

Observons d'abord que le secteur angulaire exterieur \inserted{\tikz[baseline=(n.base)]{\node[inner sep=2pt] (n) {$D_s\cap\mathrm{ext}\,\Lambda$}; \draw (n.west) .. controls (n.north west) and (n.north east) .. (n.east) .. controls (n.south east) and (n.south west) .. cycle;}} determiné par $\Lambda$ vérifie
\[
D_s\cap\mathrm{ext}\,\Lambda\subset D_s\cap\mathrm{ext}\,\Lambda^{1}\quad(\text{resp. }D_s\cap\mathrm{ext}\,\Lambda^{2}).
\]
En effet tout point $x$ de $D_s\cap\mathrm{ext}\,\Lambda$ peut être relié à un point $y$ de $\mathbb{R}^{2}$ tel que
\[
y\notin(\mathrm{int}\,\Lambda^{1}\cup\mathrm{int}\,\Lambda^{2}\cup\mathrm{int}\,\Lambda)\cup(\Lambda^{1}\cup\Lambda^{2}\cup\Lambda)
\]
\deleted{puisque} (on peut toujours trouver un tel $y$), par un arc de polygone $\Lambda_{xy}$ avec $\Lambda_{xy}\cap\Gamma=\varnothing$ \deleted{et} tel que $\Lambda_{xy}\subset\mathrm{ext}\,\Lambda$. Il est clair que $\Lambda_{xy}\cap\Lambda^{i}=\varnothing$ puisque $\Lambda_{xy}\cap{]s,s'[}=\varnothing$ car $ {]s,s'[}\deleted{\unsure{\not\subset}}\subset\mathrm{int}\,\Lambda^{i}\deleted{\unsure{\ne\varnothing|=\varnothing}}$, comme $y\notin\mathrm{int}\,\Lambda^{i}$, $i=1,2$, cela montre que $x\in\mathrm{ext}\,\Lambda^{i}$, $i=1,2$.

Cela permet de prouver que $(\Lambda^{2}\deleted{\unsure{\cap}}\mathbin{\dot{-}}{]s,s'[})\subset\overline{\mathrm{ext}\,\Lambda^{1}}$, mais le fait que $\Lambda^{2}-{]s,s'[}\subset\mathbb{R}^{2}-\Lambda^{1}$ entraîne que $\Lambda^{2}-{]s,s'[}\subset\mathrm{ext}\,\Lambda^{1}$. On montre aussi \inserted{avec le même argument} que $\Lambda^{1}-{]s,s'[}\subset\mathrm{ext}\,\Lambda^{1}$, c.q.f.d.

\uline{Lemme} $\underline{\overline{\mathrm{XIX}}}$. $\mathrm{int}\,\Lambda^{1}\subset\mathrm{ext}\,\Lambda^{2}$, $\mathrm{ext}\,\Lambda^{1}\supset\mathrm{int}\,\Lambda^{2}$.

\uline{Demonstration} : on montre que tout point de $\mathrm{int}\,\Lambda^{1}$ peut être relié\deleted{\unsure{r|e}} à un point de $\mathrm{ext}\,\Lambda^{2}$ par un arc de polygone qui ne coupe pas $\Lambda^{2}$. Comme $\Lambda^{\unsure{4|1|+}}-\{s,s'\}\subset\mathrm{ext}\,\Lambda^{2}$ soit $x\in\Lambda^{1}-\{s,s'\}$ alors \inserted{il \unsure{exìste|existe}} $D_x\subset\mathrm{ext}\,\Lambda^{2}$ satisfaisant b) de la def. 1. par rapport a $\Lambda^{1}$, donc $D_x\cap\mathrm{int}\,\Lambda^{1}\ne\varnothing$. Soi\inserted{en}t $z\in D_x\cap\mathrm{int}\,\Lambda^{1}$ et $y\in\mathrm{int}\,\Lambda^{1}$ il \unsure{exìste|existe} un arc de polygone $\Lambda_{zy}\subset\mathrm{int}\,\Lambda^{1}$ réliant $z\,(\in\mathrm{ext}\,\Lambda^{2})$ à $y$.
% archive sheet number: 44
